% Generated draft: requires mathematical and source review.
\documentclass[11pt]{article}
\usepackage[margin=25mm]{geometry}
\usepackage{fontspec}
\setmainfont{Latin Modern Roman}
\usepackage{amsmath,amssymb,mathtools,unicode-math}
\setmathfont{Latin Modern Math}
\usepackage{enumitem,needspace,xurl,hyperref}
\setlength{\parindent}{0pt}
\setlength{\parskip}{5pt}
\setlength{\emergencystretch}{3em}
\newcommand{\N}{\mathbb{N}}
\newcommand{\Z}{\mathbb{Z}}
\newcommand{\Q}{\mathbb{Q}}
\newcommand{\R}{\mathbb{R}}
\newcommand{\C}{\mathbb{C}}
\newcommand{\F}{\mathbb{F}}
\newcommand{\A}{\mathbb{A}}
\newcommand{\PP}{\mathbb{P}}
\newcommand{\E}{\mathbb{E}}
\newcommand{\eps}{\varepsilon}
\newcommand{\dd}{\,\mathrm{d}}
\begin{document}
\section*{UnsolvedMath ID: 30001854}
\textbf{Source problem code:} catalogue entry
\subsection{Definitions and mathematical statement}
Let \(I_\lambda(P;f_1,\ldots,f_m)\) denote the multilinear oscillatory integral referred to in the supplied record, with polynomial phase \(P\), input functions \(f_1,\ldots,f_m\), and oscillation parameter \(\lambda\). This notation is schematic: the record does not specify the integral, its underlying domain or measure, the maps through which the inputs enter, any amplitude or support conditions, or the admissible input spaces and normalizations. The phase is assumed to be nondegenerate in the source's intended sense, but that sense and its quantitative hypotheses are not supplied.

The supplied discussion describes a reduction by repeated Cauchy--Schwarz inequalities to a bound involving a high-order Gowers norm of one input. Write \(U^s\) for that intended norm; neither the order \(s\), the underlying group or domain, nor the normalization of this norm is specified. When the resulting estimate does not directly yield the expected decay, the discussion proposes using the Green--Tao--Ziegler inverse theorem to obtain algebraic structure for that input. The applicable version of the inverse theorem, its hypotheses, and the precise structural conclusion are not identified.

The question is whether the structural information supplied by this inverse theorem suffices to prove the desired power decay of the multilinear oscillatory integral. Schematically, power decay would mean an estimate of the form
\[
\bigl|I_\lambda(P;f_1,\ldots,f_m)\bigr|\leq C|\lambda|^{-\varepsilon},
\qquad \varepsilon>0,
\]
for sufficiently large \(|\lambda|\), with \(C\) independent of \(\lambda\). This display only explains the requested type of decay; the admissible inputs, normalization factors, dependence of \(C\) and \(\varepsilon\), and required uniformity remain unspecified. In particular, the record does not determine whether the exponent is uniform over phases or inputs.

Thus the preserved target is a question about whether the inverse-theorem structure completes the described decay argument. The supplied record does not determine a fully quantified conjecture or a standalone theorem statement without recovering the missing setup from the report.
\subsection{Short English statement}
Does the algebraic structure obtained from the Green--Tao--Ziegler inverse theorem, after a Cauchy--Schwarz reduction to a Gowers norm, suffice to establish power decay for the intended multilinear oscillatory integral with nondegenerate polynomial phase? The integral, nondegeneracy condition, inverse-theorem setting, and uniformity requirements are missing, so the question remains schematic and requires human mathematical review.
\subsection{Sources}
\begin{itemize}
\item Local UnsolvedMath registry, permanent ID \(30001854\), title: \emph{Gowers Inverse Theorems for Oscillatory Integral Decay}. Collection: Oberwolfach Reports Open Problems. Exact source problem code: \texttt{catalogue entry}. The supplied proposed year is 2011; no individual proposer is identified.
\item Supplied report citation: \emph{Real Analysis, Harmonic Analysis and Applications} (2012), Oberwolfach Reports, DOI \texttt{10.4171/owr/2011/36}, pp.\ 2045--2111. Supplied source URL: \url{https://doi.org/10.4171/owr/2011/36}. The background describes Oberwolfach Reports as documenting mathematical workshops and their open-problem discussions.
\item Supplied literature source, labeled ``OWR report 36'': \url{https://doi.org/10.4171/owr/2011/36}. The literature assessment was checked on \texttt{2026-08-21} according to the supplied record. This is supplied dated curation, not independently verified here. It labels the question open and characterizes it as an extracted follow-up question from a proof discussion, for which no definitive resolution was found. The supplied background requests independent verification of time-sensitive status and cited claims.
\item The registry assigns the statement status \texttt{corrected\_verified} and reports that the official report was checked and that the Cauchy--Schwarz reduction, Gowers-norm alternative, and inverse-theorem question were restored. These are supplied verification claims, not independent verification in this draft.
\item The source attributes the inverse theorem to Green--Tao--Ziegler. No separate bibliographic citation or theorem number for that result is supplied.
\end{itemize}
\end{document}
